% fancy-pdf.tex — mdOffice Fancy PDF style
% The document counterpart of fancy-beamer.tex: same boxes, focus boxes and
% macros, but with its own color palette (tuned for paper rather than a
% projector) and plain headings by default.
%
% ---- FRONTMATTER EXAMPLE ---------------------------------------------------
%
% ---
% title:   "Document Title"
% author:  "Your Name"
% date:    "Date"
%
% mdoffice:
%   make-pdf:  true
%   doc-style: "fancy-pdf"
% ---
%
% ---- HEADING COLORS --------------------------------------------------------
% Headings are plain (no background) by default. To give ONE heading a
% colored bar, put a macro on the line right after it (same as fancy-beamer):
%
%   ## My Heading
%   \redheader
%
% Available: \myheader \blueheader \redheader \greenheader \cyanheader \grayheader
% The color applies to that heading only; the next heading is plain again.
% (The pdf_header_color.lua filter moves the macro in front of the heading,
% so writing it before the heading works too.)
%
% ---- THEOREM BOXES ---------------------------------------------------------
% Numbered, with title and cross-reference label:
%   \begin{bluebox}{Title}{label} ... \end{bluebox}
%   \begin{redbox}{Title}{label}  ... \end{redbox}
%   \begin{greenbox}{Title}{label}... \end{greenbox}
%   \begin{cyanbox}{Title}{label} ... \end{cyanbox}
%   \begin{graybox}{Title}{label} ... \end{graybox}
% Starred form (no counter):
%   \begin{bluebox*}{Title} ... \end{bluebox*}
%
% ---- FOCUS BOXES -----------------------------------------------------------
% No title, lightly tinted background — for key formulas, remarks, warnings:
%   \begin{myfocus}   ... \end{myfocus}
%   \begin{bluefocus} ... \end{bluefocus}
%   \begin{redfocus}  ... \end{redfocus}
%   \begin{greenfocus}... \end{greenfocus}
%   \begin{cyanfocus} ... \end{cyanfocus}
%   \begin{grayfocus} ... \end{grayfocus}
%
% ---- WARNING / INFO MACROS -------------------------------------------------
%   \warn                              — red triangle + "Not part of the syllabus"
%   \warnCustomMsg[color]{text}        — custom triangle message
%   \info                              — blue circle  + "For your information"
%   \infoCustomMsg[color]{text}        — custom info  message
%
% ---- MATH MACROS -----------------------------------------------------------
%   \norm{v}   →  ‖v‖        \abs{x}   →  |x|
% ----------------------------------------------------------------------------


% ---------------------------------------------------------------
% PACKAGES
\usepackage{geometry}
\usepackage{graphicx}
\usepackage{float}
\floatplacement{figure}{H}
\usepackage{placeins}
\usepackage{array}
\usepackage{xcolor}
\usepackage[most]{tcolorbox}    % theorem/focus boxes
\usepackage{fancyhdr}
\usepackage{lastpage}
\usepackage{titling}            % \thetitle / \theauthor / \thedate
\usepackage{etoolbox}
\usepackage[explicit]{titlesec} % heading styles (explicit: title passed as #1)
\usepackage{fontawesome5}       % icons for \warn / \info
\usepackage{xspace}             % smart spacing after macros
\usepackage{icomma}             % no extra space after "," in math mode, e.g. $2,0$
\usepackage{colortbl}           % \rowcolor / \cellcolor for table shading
% ---------------------------------------------------------------

% Allow @ in internal macro names (\fancypdf@...) for the rest of this file.
\makeatletter


% ---------------------------------------------------------------
% COLORS — fancy-pdf's own palette (fancy-beamer.tex has a separate one).
% Paper shows colors darker and less luminous than a projector, so the
% base colors are slightly deeper and the focus tints slightly lighter.
%
% 1. Base colors — raw RGB palette, used for inline text and as source for derived colors
\definecolor{myblue}{RGB}{25, 80, 150}
\definecolor{myred}{RGB}{170, 35, 35}
\definecolor{mygreen}{RGB}{40, 115, 45}
\definecolor{mycyan}{RGB}{0, 140, 150}
\definecolor{mygray}{RGB}{70, 70, 70}
%
% 2. Heading colors — for heading bars and theorem box frames
%    Adjust the ! expression to darken/shift without touching the base RGB
\colorlet{myblueheading}{myblue!80}
\colorlet{myredheading}{myred!80}
\colorlet{mygreenheading}{mygreen!80}
\colorlet{mycyanheading}{mycyan!80}
\colorlet{mygrayheading}{mygray!80}
%
% 3. Focus colors — light tints of the base colors, for focus box backgrounds
\colorlet{mybluefocus}{myblue!8}
\colorlet{myredfocus}{myred!8}
\colorlet{mygreenfocus}{mygreen!8}
\colorlet{mycyanfocus}{mycyan!10}
\colorlet{mygrayfocus}{mygray!6}
%
% 4. Plain heading text color (headings without a \xxxheader macro)
\colorlet{myplainheading}{black}
%
% Header color names used by \blueheader etc.
\def\HeaderDefault{myblueheading}
\def\HeaderBlue{myblueheading}
\def\HeaderRed{myredheading}
\def\HeaderGreen{mygreenheading}
\def\HeaderCyan{mycyanheading}
\def\HeaderGray{mygrayheading}
% ---------------------------------------------------------------


% ---------------------------------------------------------------
% PAGE LAYOUT, HEADER AND FOOTER (same contract as standard-pdf.tex)
\ifdefined\logopath
  \geometry{top=0.2cm, bottom=2.54cm, left=2.54cm, right=2.54cm,
            headheight=2.4cm, headsep=0.6cm, includehead}
\else
  \geometry{top=2.00cm, bottom=2.00cm, left=2.5cm, right=2.5cm}
\fi

\newcommand{\mdpagefooter}{%
  \ifdefined\hidepagenumbers\else
    \fancyfoot[R]{\thepage}%
  \fi
}

% Left slot: logo and \leftheader side by side, logo first.
\newcommand{\mdpageheader}{%
  \fancyhf{}%
  \fancyhead[L]{%
    \ifdefined\logopath\includegraphics[height=\logoheight]{\logopath}\ifdefined\leftheader\quad\fi\fi%
    \ifdefined\leftheader{\fontsize{9}{8}\selectfont\leftheader}\fi}%
  \ifdefined\rightheader
    \fancyhead[R]{{\fontsize{9}{8}\selectfont\rightheader}}%
  \fi
  \renewcommand{\headrulewidth}{0pt}%
}

\pagestyle{fancy}
\mdpageheader
\mdpagefooter

\fancypagestyle{plain}{\mdpageheader\mdpagefooter}
\fancypagestyle{titlepage}{\mdpageheader}
% ---------------------------------------------------------------


% ---------------------------------------------------------------
% TITLE PAGE — flush-left, own page, title in the theme's blue
\def\fancypdf@empty{}
\renewcommand{\maketitle}{%
  \ifdefined\hidefirstpagenumber\thispagestyle{titlepage}\else\thispagestyle{fancy}\fi%
  \noindent{\color{myblueheading}\Huge\bfseries\thetitle\par}%
  \vspace{0.4em}%
  \ifdefined\mdsubtitle{\noindent\Large\color{mygray}\mdsubtitle\par}\fi%
  \vspace{0.8em}%
  {\noindent\color{myblueheading}\rule{\linewidth}{1.2pt}\par}%
  \vspace{0.6em}%
  \ifx\theauthor\fancypdf@empty\else\noindent\theauthor\par\fi%
  \ifx\thedate\fancypdf@empty\else\noindent\thedate\par\fi%
  \newpage%
  \ifdefined\hidefirstpagenumber\setcounter{page}{1}\fi%
}

\ifdefined\hidefirstpagenumber
  \pretocmd{\tableofcontents}{\pagenumbering{roman}}{}{}
  \apptocmd{\tableofcontents}{\newpage\pagenumbering{arabic}}{}{}
\fi
% ---------------------------------------------------------------


% ---------------------------------------------------------------
% HEADING COLOR SYSTEM
%
% \fancypdf@headcolor holds the color for the NEXT heading only; empty
% means plain. \blueheader etc. set it, and every heading resets it after
% use, so colored bars are always opt-in. \gdef keeps it global across
% groups (e.g. when the macro ends up inside a pandoc div).
%
% \fancypdf@num is filled from titlesec's label slot, which titlesec only
% typesets for numbered headings — so the number appears inside the bar
% when `numbersections: true` and is absent otherwise.
\gdef\fancypdf@headcolor{}
\gdef\fancypdf@num{}

\newcommand{\blueheader}{\gdef\fancypdf@headcolor{\HeaderBlue}}
\newcommand{\redheader}{\gdef\fancypdf@headcolor{\HeaderRed}}
\newcommand{\greenheader}{\gdef\fancypdf@headcolor{\HeaderGreen}}
\newcommand{\cyanheader}{\gdef\fancypdf@headcolor{\HeaderCyan}}
\newcommand{\grayheader}{\gdef\fancypdf@headcolor{\HeaderGray}}
\newcommand{\myheader}{\gdef\fancypdf@headcolor{\HeaderDefault}}

% #1 = font size of the heading level, #2 = heading text. Draws a
% full-width colored bar with white text when a color is pending,
% otherwise a plain heading in \myplainheading.
% The bar uses tcolorbox's default arc, so its rounded corners match the
% tinted area of the focus boxes. Change `arc` here for rounder/sharper bars.
% tcolorbox resets the font inside the box, so the size is passed in again.
\newcommand{\fancypdf@heading}[2]{%
  \ifx\fancypdf@headcolor\fancypdf@empty
    {\color{myplainheading}\fancypdf@num#2}%
  \else
    \edef\fancypdf@bgcol{\fancypdf@headcolor}%
    \tcbox[
        colback=\fancypdf@bgcol, colframe=\fancypdf@bgcol, coltext=white,
        fontupper=#1\bfseries,
        boxrule=0pt, arc=1mm, nobeforeafter, tcbox raise base,
        boxsep=0pt, left=4pt, right=4pt, top=3pt, bottom=2pt]{%
      \parbox{\dimexpr\linewidth-8pt\relax}{\raggedright\strut\fancypdf@num#2\strut}}%
  \fi
  \gdef\fancypdf@headcolor{}%
  \gdef\fancypdf@num{}%
}

\titleformat{\section}[block]{\normalfont\Large\bfseries}%
  {\xdef\fancypdf@num{\thesection\quad}}{0pt}{\fancypdf@heading{\Large}{#1}}
\titleformat{\subsection}[block]{\normalfont\large\bfseries}%
  {\xdef\fancypdf@num{\thesubsection\quad}}{0pt}{\fancypdf@heading{\large}{#1}}
\titleformat{\subsubsection}[block]{\normalfont\normalsize\bfseries}%
  {\xdef\fancypdf@num{\thesubsubsection\quad}}{0pt}{\fancypdf@heading{\normalsize}{#1}}
% ---------------------------------------------------------------


% ---------------------------------------------------------------
% TCOLORBOX GLOBAL DEFAULTS
\tcbset{
  colback=white,
  fonttitle=\bfseries,
  breakable,
  enhanced,
  top=2pt, bottom=2pt, left=2pt, right=2pt,
  boxsep=3pt
}

% Theorem-style boxes: \begin{bluebox}{Title}{label} ... \end{bluebox}
% All share a single counter so numbering is consistent across colors.
\newtcbtheorem{bluebox}{}%
  {colframe=myblueheading}{bluebox}

\newtcbtheorem[use counter from=bluebox]{redbox}{}%
  {colframe=myredheading}{redbox}

\newtcbtheorem[use counter from=bluebox]{greenbox}{}%
  {colframe=mygreenheading}{greenbox}

\newtcbtheorem[use counter from=bluebox]{cyanbox}{}%
  {colframe=mycyanheading}{cyanbox}

\newtcbtheorem[use counter from=bluebox]{graybox}{}%
  {colframe=mygrayheading}{graybox}

% Focus boxes: no title, lightly tinted background
\newtcolorbox{myfocus}{colback=mybluefocus,   colframe=white}
\newtcolorbox{bluefocus}{colback=mybluefocus, colframe=white}
\newtcolorbox{redfocus}{colback=myredfocus,   colframe=white}
\newtcolorbox{greenfocus}{colback=mygreenfocus, colframe=white}
\newtcolorbox{cyanfocus}{colback=mycyanfocus, colframe=white}
\newtcolorbox{grayfocus}{colback=mygrayfocus, colframe=white}
% ---------------------------------------------------------------


% ---------------------------------------------------------------
% WARNING AND INFO MACROS
\newcommand{\warn}{%
  \textcolor{myred}{\faExclamationTriangle\ \textbf{Not part of the syllabus}}\xspace%
}
\newcommand{\warnCustomMsg}[2][myred]{%
  \textcolor{#1}{\faExclamationTriangle\ \textbf{#2}}\xspace%
}
\newcommand{\info}{%
  \textcolor{myblue}{\faInfoCircle\ \textbf{For your information}}\xspace%
}
\newcommand{\infoCustomMsg}[2][myblue]{%
  \textcolor{#1}{\faInfoCircle\ \textbf{#2}}\xspace%
}
% ---------------------------------------------------------------


% ---------------------------------------------------------------
% MATH MACROS
\providecommand{\norm}[1]{\left\lVert#1\right\rVert}
\providecommand{\abs}[1]{\left\lvert#1\right\rvert}
% ---------------------------------------------------------------

\makeatother
